% ch09 译文：原文 §5d 限制到子群概形、§5e 群概形作用的乘积（main.tex 第 2671--3290 行）
\subsection{限制到子群概形}

设 $G$ 为群概形，$H\subseteq G$ 为闭子群概形，包含态射记为 $i_H : H \hookrightarrow G$。概形 $X$ 上的 $G$-作用
$\sigma_G : G \times X \rightarrow X$ 自然地诱导出 $X$ 上的 $H$-作用，由
$$
\sigma_H
:=
(i_H \times \id_X)
\comp
\sigma
:
H\times X\rightarrow X.
$$
给出。

\begin{defn}[等变层的限制]\label{defn:RestrictionOfEquivariantSheaves}
对任意 $G$-等变 $\mathcal{O}_X$-模 $(\mathcal{F}, \theta) \in \Mod_G(X)$，其到 $H$ 的\emph{限制}定义为 $H$-等变 $\mathcal{O}_X$-模
$$
(\mathcal{F}, \theta)|_H := (\mathcal{F}, \theta|_H),
$$
其中受限的 $H$-线性化由 $G$-线性化拉回得到
$$
\theta|_H 
:= 
(i_H \times \id_X)^* \theta 
: 
\sigma_H^* \mathcal{F} 
\xrightarrow{\cong} 
\left(\pr_X^{H \times X}\right)^* \mathcal{F}
$$
这一对象层面的操作连同态射上的恒等映射，自然地定义出一个函子
$$
\Res_H^G : \Mod_G(X) \rightarrow \Mod_H(X),
$$
称为\emph{限制函子}。它也在相应的全子范畴上诱导限制函子：
$$
\QCoh_G(X) \rightarrow \QCoh_H(X),
\quad
\Coh_G(X) \rightarrow \Coh_H(X),
\quad\text{与}\quad
\VB_G(X) \rightarrow \VB_H(X).
$$
\end{defn}

\begin{yizhu}[限制的方向]
限制只换作用群，不动层本身：$\mathcal{F}$ 保持不变，改变的只是线性化 $\theta$。注意 $\theta|_H$ 沿 $i_H\times\id_X$ 拉回，同构两端都落在 $H\times X$ 上；若把 $i_H$ 与 $\sigma_G$ 的复合顺序颠倒，两端就不在同一个空间上，限制便无从谈起。这是本节最容易搞反的一步。
\end{yizhu}

下面这个引理是标准的：

\begin{lemma}\cite[Exercise 2.8]{Waterhouse}
\label{lem:DiagonalizableSurjective}
设 $G$ 是域 $\field$ 上的对角化群概形，$H \subseteq G$ 为闭子群概形。则 $H$ 也是对角化的，且特征群上的限制映射
$
i_H^* : \widehat{G} \rightarrow \widehat{H}
$
是满射，其核为 $H^\perp$，其中
$$
H^\perp
:=
\{ \cha \in \widehat{G} \mid \cha(h) = 1 \text{ 对一切 } h \in H \}
\subseteq
\widehat{G}
$$
称为 $H$ 在 $\widehat{G}$ 中的零化子。
\end{lemma}

\begin{proof}
由于 $H$ 是对角化群概形 $G\cong \Spec(\field[\widehat{G}])$ 的闭子群概形，存在 $\field[\widehat{G}]$ 的 Hopf 理想 $I$，使得 $H\cong\Spec(\field[\widehat{G}]/I)$ 且 $\mathcal{O}_H(H)\cong\field[\widehat{G}]/I$。

我们断言
$$
I=\left<\cha-1|\cha\in H^\perp\right>
\subseteq
\field[\widehat{G}].
$$

对任意 $\cha \in H^\perp\subseteq\field[\widehat{G}]\cong\mathcal{O}_G(G)$，有 $\cha|_H=1_{\mathcal{O}_H(H)}\in\mathcal{O}_H(H)$，从而 $(\cha-1)|_H=0\in \mathcal{O}_H(H)$。这说明 $\langle \cha - 1 \mid \cha \in H^\perp \rangle \subseteq I$。

反之，设 $f = \sum_{\cha \in \widehat{G}} c_{\cha} \cha \in I\subseteq\field[\widehat{G}]\cong\mathcal{O}_G(G)$。按各项在 $\mathcal{O}_H(H)$ 中的限制像归并，对 $f|_H$ 求值：
$$
0 = f|_H = \sum_{\mu} \left( \sum_{\cha|_H = \mu} c_{\cha} \right) \mu \quad \text{在 } \mathcal{O}_H(H) \text{ 中},
$$
其中 $\mu$ 取遍特征在 $\mathcal{O}_H(H)$ 中的互不相同的非零像。由于每个 $\mu$ 都是 Hopf 代数 $\mathcal{O}_H(H)$ 中的群像元素，这些互不相同的元素在 $\field$ 上线性无关
\cite[\S2.2]{Waterhouse}。于是每个 $\mu$ 的系数都必须为零，即对所有这样的 $\mu$ 都有 $\sum_{\cha|_H = \mu} c_{\cha} = 0$。

% [原文 p.38]
对求和中出现的每个 $\mu$，固定一个代表元 $\cha_0 \in \widehat{G}$，使得 $\cha_0|_H = \mu$。任何其他限制到 $\mu$ 的特征 $\cha$ 都满足 $(\cha \cha_0^{-1})|_H = 1$，即 $\cha \cha_0^{-1} \in H^\perp$。于是可写 $\cha = h\cha_0$，其中 $h \in H^\perp$。这样就能把对应于 $\mu$ 的那部分和改写为
$$
\sum_{h \in H^\perp} c_{h\cha_0} (h\cha_0) = \sum_{h \in H^\perp} c_{h\cha_0} (h - 1)\cha_0 + \left( \sum_{h \in H^\perp} c_{h\cha_0} \right) \cha_0.
$$
由于系数之和 $\sum c_{h\cha_0}$ 恰为零，第二项消失，只剩下形如 $h-1$（$h \in H^\perp$）的元素的线性组合。对所有这样的 $\mu$ 求和，即得 $f \in \langle \cha - 1 \mid \cha \in H^\perp \rangle$。

于是 $I = \langle \cha - 1 \mid \cha \in H^\perp \rangle$。由环的第一同构定理得到自然同构 $\mathcal{O}_H(H) \cong \field[\widehat{G}]/I \cong \field[\widehat{G}/H^\perp]$，故 $H$ 是对角化的，且 $\widehat{G}/H^\perp\cong\widehat{H}$。
\end{proof}

\begin{yizhu}[零化子零化了什么]
$H^\perp$ 由「在 $H$ 上恒取 $1$」的特征组成，也就是 $H$ 看不见的那些权。因此 $H$-等变层的权分解并不直接按 $\widehat{G}$ 铺开：每个 $H^\perp$-陪集里的 $\widehat{G}$-权会被并成同一个 $\widehat{H}$-权。推论 \ref{crl:orbit-category-GH-restriction} 里按 $H^\perp$ 取商，正是这一步的范畴化。
\end{yizhu}

\begin{prop}[同态空间的 $H$-不变量]
\label{prop:HInvariantsOfHomSpace}
设 $G$ 是对角化群概形、作用在概形 $X$ 上，$H\subseteq G$ 为闭子群概形。对任意
$(\mathcal{F}, \theta)$、
$(\mathcal{F}', \theta') \in \Mod_G(X)$，只要 $\mathcal{F}$ 有限呈现，$\Mod_H(X)$ 中从 $(\mathcal{F}, \theta|_H)$ 到 $(\mathcal{F}', \theta'|_H)$ 的同态空间就有典范分解
$$
\operatorname{Hom}_{\Mod_H(X)}\left( (\mathcal{F}, \theta|_H), (\mathcal{F}', \theta'|_H) \right)
=
\bigoplus_{\cha \in H^\perp} \operatorname{Hom}_{\Mod_G(X)}\left( (\mathcal{F}, \theta)(\cha), (\mathcal{F}', \theta') \right).
$$
\end{prop}

\begin{proof}
由引理 \ref{lem:DiagonalizableSurjective}，$H$ 是对角化的。因此由命题 \ref{prop:WeightSpaceDecompositionOfHomSpace} 可得两个分解
\begin{align*}
\operatorname{Hom}_{\Mod(X)}(\mathcal{F}, \mathcal{F}') 
&=
\bigoplus_{\cha \in \widehat{G}} \operatorname{Hom}_{\Mod_G(X)}\left( (\mathcal{F}, \theta)(\cha), (\mathcal{F}', \theta') \right)
\\
&=
\bigoplus_{\chaa \in \widehat{H}} \operatorname{Hom}_{\Mod_H(X)}\left( (\mathcal{F}, \theta|_H)(\chaa), (\mathcal{F}', \theta'|_H) \right).
\end{align*}

我们断言对每个 $\cha\in\widehat{G}$ 都有
$$
\operatorname{Hom}_{\Mod_G(X)}\left( (\mathcal{F}, \theta)(\cha), (\mathcal{F}', \theta') \right)
\subseteq
\operatorname{Hom}_{\Mod_H(X)}\left( (\mathcal{F}, \theta|_H)(i_H^*\cha), (\mathcal{F}', \theta'|_H) \right).
$$
事实上，$f\in\operatorname{Hom}_{\Mod_G(X)}\left( (\mathcal{F}, \theta)(\cha), (\mathcal{F}', \theta') \right)$ 满足
$$
\sigma_G^* f\comp \theta'
=
\left(\left(\pr_G^{G\times X}\right)^*\cha \cdot \theta\right)
\comp
\left(\pr_X^{G\times X}\right)^* f
$$
% [原文 p.39]
在两侧取 $(i_H^*\times \id_X)^*$，得到
$$
\sigma_H^*f\comp\left(\theta'|_H\right)
=
\left(
\left(\pr_H^{H\times X}\right)^*
i_H^*\cha
\cdot
\theta|_H
\right)
\comp
\left(\pr_X^{H\times X}\right)^*f,
$$
其中
$
\theta|_H
(i_H^*\cha)
=
\left(
\left(\pr_H^{H\times X}\right)^*
i_H^*\cha
\cdot
\theta|_H
\right).
$
这说明
$
f\in
\operatorname{Hom}_{\Mod_H(X)}\left( (\mathcal{F}, \theta|_H)(i_H^*\cha), (\mathcal{F}', \theta'|_H) \right),
$
断言得证。

因此，对每个 $\chaa\in\widehat{H}$ 都有如下子分解
$$
\operatorname{Hom}_{\Mod_H(X)}\left( (\mathcal{F}, \theta|_H)(\chaa), (\mathcal{F}', \theta'|_H) \right)
=
\bigoplus_{\substack{\cha \in \widehat{G} \\[1mm] i_H^*\cha = \chaa}} \operatorname{Hom}_{\Mod_G(X)}\left( (\mathcal{F}, \theta)(\cha), (\mathcal{F}', \theta') \right).
$$
命题即为 $\chaa=1$ 的特殊情形。
\end{proof}

与推论 \ref{thm:EmbeddingOfTheOrbitCategory} 类似，存在把 $G$-等变向量丛限制为 $H$-等变向量丛的轨道范畴嵌入。

\begin{crl}\label{crl:orbit-category-GH-restriction}
    设 $G$ 为群概形，$H\subseteq G$ 为 $G$ 的子群概形。则存在阿贝尔范畴的嵌入 \[
        \QCoh_G (X)/H^{\perp}\hookrightarrow\QCoh_H (X).
    \]
\end{crl}
\begin{proof}
按定义，轨道范畴中的同态空间为
$$
\operatorname{Hom}_{\QCoh_G(X)/H^\perp}\left( (\mathcal{F}, \theta), (\mathcal{F}', \theta') \right)
=
\bigoplus_{\cha \in H^\perp} \operatorname{Hom}_{\QCoh_G(X)}\left( (\mathcal{F}, \theta)(\cha), (\mathcal{F}', \theta') \right).
$$
由命题 \ref{prop:HInvariantsOfHomSpace}，这一直和恰好等于 $\QCoh_H(X)\subseteq\Mod_H(X)$ 中的同态空间
\[\operatorname{Hom}_{\QCoh_H(X)}
\left((\mathcal{F}, \theta|_H), (\mathcal{F}', \theta'|_H) \right)\]。
因此，在对象上取 $(\mathcal{F}, \theta) \mapsto (\mathcal{F},\theta|_H)$、在态射上取恒等映射所定义的函子 $\QCoh_G(X)/H^\perp \hookrightarrow\QCoh_H(X)$ 是全忠实的。
\end{proof}

\begin{yizhu}[推论略去的前提]
推论的证明直接调用命题 \ref{prop:HInvariantsOfHomSpace}，而后者要求 $\mathcal{F}$ 有限呈现，推论却对一般拟凝聚层陈述。原文一句带过：通常的补法是利用拟凝聚层可写成有限呈现对象的滤过余极限，而两侧的同态空间都与该余极限交换。补齐这一步并不困难，但不该把该命题读成对任意 $\mathcal{F}$ 无条件成立。
\end{yizhu}

\subsection{群概形作用的乘积}

设群概形 $G$ 与 $H$ 分别作用在概形 $X$ 与 $Y$ 上，作用态射分别为
$\sigma_X : G \times X \rightarrow X$
与
$\sigma_Y : H \times Y \rightarrow Y$。
它们自然地诱导乘积群概形 $G\times H$ 在乘积概形 $X\times Y$ 上的作用，其作用态射为
$$
\sigma_{X\times Y}
:=
\sigma_X\times \sigma_Y
:
(G\times H)\times (X\times Y)
\rightarrow
X\times Y.
$$
\begin{defn}[等变层的外张量积]
\label{defn:ExternalTensorProductOfEquivariantModules}
设 $(\mathcal{F},\theta)\in\Mod_G(X)$ 与
$(\mathcal{G},\eta)\in\Mod_H(Y)$。
$(\mathcal{F},\theta)$ 与 $(\mathcal{G},\eta)$ 的\emph{外张量积}
$$
(\mathcal{F},\theta)\boxtimes(\mathcal{G},\eta)
:=
\left(
\mathcal{F}\boxtimes\mathcal{G},
\
\theta\boxtimes\eta
\right)
\in\Mod_{G\times H}(X\times Y)
$$
定义为 $\mathcal{O}_{X\times Y}$-模
$
\mathcal{F}\boxtimes\mathcal{G}
:=
\pi_X^*\mathcal{F}
\underset{\mathcal{O}_{X\times Y}}{\otimes}
\pi_Y^*\mathcal{G}
$
% [原文 p.40]
连同 $G\times H$-线性化
$$
\theta\boxtimes\eta
:=
\tilde{\pi}_{G\times X}^*\theta\otimes\tilde{\pi}_{H\times Y}^*\eta
:
\sigma_{X\times Y}^*(\mathcal{F}\boxtimes\mathcal{G})
\xrightarrow{\cong}
\left(\pr_{X\times Y}^{(G\times H)\times(X\times Y)}\right)^*
(\mathcal{F}\boxtimes\mathcal{G})
$$
其中
\begin{align*}
\tilde{\pi}_{G\times X}^*\theta
&:
\tilde{\pi}_{G\times X}^*(\sigma_X^*\mathcal{F})
\xrightarrow{\cong}
\tilde{\pi}_{G\times X}^*\left(\left(\pr_X^{G\times X}\right)^{\ast}\mathcal{F}\right),
\\[2mm]
\tilde{\pi}_{H\times Y}^*\eta
&:
\tilde{\pi}_{H\times Y}^*(\sigma_Y^*\mathcal{G})
\xrightarrow{\cong}
\tilde{\pi}_{H\times Y}^*\left(\left(\pr_Y^{H\times Y}\right)^*\mathcal{G}\right)
\end{align*}
分别是 $\theta$ 与 $\eta$ 沿自然投影
\begin{align*}
\tilde{\pi}_{G\times X}
&:=
\pr_{G\times X}^{(G\times H)\times(X\times Y)}
:
(G\times H)\times(X\times Y)
\rightarrow
G\times X,
\\[2mm]
\tilde{\pi}_{H\times Y}
&:=
\pr_{H\times Y}^{(G\times H)\times(X\times Y)}
:
(G\times H)\times(X\times Y)
\rightarrow
H\times Y,
\end{align*}
的拉回，而张量积
$
\tilde{\pi}_{G\times X}^*\theta\otimes\tilde{\pi}_{H\times Y}^*\eta
$
在如下典范同构下给出
\begin{align*}
\sigma_{X\times Y}^*(\mathcal{F}\boxtimes\mathcal{G})
&=
\sigma_{X\times Y}^*
\Big(
\pi_X^*\mathcal{F}
\underset{\mathcal{O}_{X\times Y}}{\otimes}
\pi_Y^*\mathcal{G}
\Big)
\\[2mm]
&\cong
(\sigma_{X\times Y}\comp\pi_X)^*\mathcal{F}
\underset{\mathcal{O}_{(G\times H)\times(X\times Y)}}{\otimes}
(\sigma_{X\times Y}\comp\pi_Y)^*\mathcal{G}
\\[2mm]
&=
(\tilde{\pi}_{G\times X}\comp \sigma_X)^*\mathcal{F}
\underset{\mathcal{O}_{(G\times H)\times(X\times Y)}}{\otimes}
(\tilde{\pi}_{H\times Y}\comp \sigma_Y)^*\mathcal{G}
\\[2mm]
&\cong
\tilde{\pi}_{G\times X}^*\left(\sigma_X^*\mathcal{F}\right)
\underset{\mathcal{O}_{(G\times H)\times(X\times Y)}}{\otimes}
\tilde{\pi}_{H\times Y}^*\left(\sigma_Y^*\mathcal{G}\right)
\end{align*}
与
\begin{align*}
\left(\pr_{X\times Y}^{(G\times H)\times(X\times Y)}\right)^*
(\mathcal{F}\boxtimes\mathcal{G})
&=
\left(\pr_{X\times Y}^{(G\times H)\times(X\times Y)}\right)^*
\Big(
\pi_X^*\mathcal{F}
\underset{\mathcal{O}_{X\times Y}}{\otimes}
\pi_Y^*\mathcal{G}
\Big)
\\
&\cong
\left(\pr_{X\times Y}^{(G\times H)\times(X\times Y)}\comp\pi_X\right)^*\mathcal{F}
\underset{\mathcal{O}_{(G\times H)\times(X\times Y)}}{\otimes}
\left(\pr_{X\times Y}^{(G\times H)\times(X\times Y)}\comp\pi_Y\right)^*\mathcal{G}
\\
&\cong
\left(\tilde{\pi}_{G\times X}\comp\pr_X^{G\times X}\right)^*\mathcal{F}
\underset{\mathcal{O}_{(G\times H)\times(X\times Y)}}{\otimes}
\left(\tilde{\pi}_{H\times Y}\comp\pr_Y^{H\times Y}\right)^*\mathcal{G}
\\
&\cong
\tilde{\pi}_{G\times X}^*\left(\left(\pr_X^{G\times X}\right)^*\mathcal{F}\right)
\underset{\mathcal{O}_{(G\times H)\times(X\times Y)}}{\otimes}
\tilde{\pi}_{H\times Y}^*\left(\left(\pr_Y^{H\times Y}\right)^*\mathcal{G}\right).
\end{align*}

对态射
$
f
\in
\Hom_{\Mod_G(X)}
((\mathcal{F},\theta),(\mathcal{F}',\theta'))
$
与
$
g
\in
\Hom_{\Mod_H(Y)}
((\mathcal{G},\eta),(\mathcal{G}',\eta')),
$
定义
$$
f\boxtimes g
:=
\pi_X^*f
\otimes
\pi_Y^*g
\in
\Hom_{\Mod_{G\times H}(X\times Y)}
\big(
\left(
\mathcal{F}\boxtimes\mathcal{G},
\
\theta\boxtimes\eta
\right),
\left(
\mathcal{F}'\boxtimes\mathcal{G}',
\
\theta'\boxtimes\eta'
\right)
\big).
$$

这定义出一个双函子
$$
\boxtimes:
\Mod_G(X)\times\Mod_H(Y)
\rightarrow
\Mod_{G\times H}(X\times Y),
$$
它限制到等变（拟）凝聚层与等变向量丛上。
\end{defn}

\begin{conceptbox}{等变层的外张量积}
给定 $G$-等变层 $(\mathcal{F},\theta)$ 与 $H$-等变层 $(\mathcal{G},\eta)$，外张量积先在底空间上取 $\pi_X^*\mathcal{F}\otimes\pi_Y^*\mathcal{G}$，再把两个线性化沿自然投影拉回并张量起来，得到 $G\times H$-等变层。它对两个变量分别可加，并给出双函子 $\Mod_G(X)\times\Mod_H(Y)\to\Mod_{G\times H}(X\times Y)$；命题 \ref{prop:BoxTensorFunctorFullyFaithful} 进一步说明它在同态层面全忠实。本文需要它，是因为扭 Künneth 公式要把两侧的等变导出范畴「相乘」成乘积空间上的等变导出范畴，而外张量积正是这个乘法在层层面的原型。
\par\smallskip\noindent{\footnotesize 参考：nLab，\emph{equivariant sheaf}；Waterhouse, \emph{Introduction to Affine Group Schemes}, GTM 66, Springer, 1979, 第 2 章。}
\end{conceptbox}

\begin{yizhu}[外张量积为什么不必担心符号]
$G$ 只作用在 $X$ 侧、$H$ 只作用在 $Y$ 侧，两个线性化落在互不相交的变量上，因此可以逐项张量而不产生符号：这里张量的是 $\mathcal{O}$-模而不是复形，没有 Koszul 号出场的机会。真正需要核对的是两个投影 $\tilde{\pi}_{G\times X}$、$\tilde{\pi}_{H\times Y}$ 的方向，它们把 $(G\times H)\times(X\times Y)$ 分别投到 $G\times X$ 与 $H\times Y$。
\end{yizhu}

\begin{prop}\label{prop:BoxTensorFunctorFullyFaithful}
函子
$$
\boxtimes:
\Mod_G(X)\otimes\Mod_H(Y)
\hookrightarrow
\Mod_{G\times H}(X\times Y)
$$
是全忠实的，即
\begin{align*}
\boxtimes:
\Hom_{\Mod_G(X)}&
\big(
(\mathcal{F},\theta),
(\mathcal{F}',\theta')
\big)
\
\otimes
\
\Hom_{\Mod_H(Y)}
\big(
(\mathcal{G},\eta),
(\mathcal{G}',\eta')
\big)
\\ 
&\xrightarrow{\cong}
\Hom_{\Mod_{G\times H}(X\times Y)}
\big(
(\mathcal{F},\theta)\boxtimes(\mathcal{G},\eta),
(\mathcal{F}',\theta')\boxtimes(\mathcal{G}',\eta')
\big)
\end{align*}
\end{prop}

\begin{keytakeaways}
  \item 限制是本节的基本操作：把 $G$-等变层的线性化沿 $H\hookrightarrow G$ 限制，就得到 $H$-等变层，并给出限制函子 $\Res_H^G$；层本身不动，动的只是线性化。
  \item 对角化群概形使权分解可以「沿子群重排」：命题 \ref{prop:HInvariantsOfHomSpace} 把 $\Mod_H(X)$ 的同态空间写成 $\bigoplus_{\cha\in H^\perp}\operatorname{Hom}_{\Mod_G(X)}((\mathcal{F},\theta)(\cha),(\mathcal{F}',\theta'))$，于是 $H$-不变部分恰好由 $H^\perp$ 标号。
  \item 这条分解的范畴化就是推论 \ref{crl:orbit-category-GH-restriction}：轨道范畴 $\QCoh_G(X)/H^\perp$ 全忠实地嵌入 $\QCoh_H(X)$，与前面 $\QCoh_G(X)/\widehat{G}\hookrightarrow\QCoh(X)$ 的嵌入是同一套机制。
  \item 反向的拼装由外张量积负责：$\boxtimes$ 把两个等变范畴送进乘积空间的等变范畴，且在全子范畴上全忠实（命题 \ref{prop:BoxTensorFunctorFullyFaithful}）。这正是下一节扭 Künneth 公式所需的乘法。
\end{keytakeaways}
